%======================================================================
\section{Introduction}
\label{sec:intro}
%======================================================================

The dynamics of gluon saturation \cite{Morreale:2021pnn,Albacete:2014fwa,Blaizot:2016qgz,Mueller:1985wy}
at small-$x$ can be systematically studied within the framework of the Color Glass Condensate
(CGC) effective theory \cite{Gelis:2010nm,Weigert:2005us,McLerran:1993ni,McLerran:1993ka,Balitsky:2001gj,Garcia-Montero:2025hys}.
In this regime the rapidity evolution of the gluon distributions is governed by the nonlinear
BK and JIMWLK equations \cite{Iancu:2001ad,Iancu:2000hn,Weigert:2000gi,Kovner:2000pt,Jalilian-Marian:1997jx,Jalilian-Marian:1997gr,Jalilian-Marian:1996mkd,Kovchegov:1999yj,Kovchegov:1999ua,Lublinsky:2016meo,Kovner:2014lca,Kovner:2013ona}.
While these equations describe the evolution of Wilson line correlators, physical observables
such as single and double inclusive gluon production \cite{Kovner:2011pe,Altinoluk:2018byz} and,
more generally, particle production \cite{Iancu:2020mos,Dominguez:2012ad,Altinoluk:2011qy,Blaizot:2004wv,Kovchegov:2006qn,Dumitru:2005gt}
are obtained from matrix elements of operators evaluated between boosted hadronic states.
An explicit and controlled construction of the light-cone wave function (LCWF) of a fast-moving
hadron is therefore of both theoretical and phenomenological interest.

From the light-cone perspective, boosting a hadron by a rapidity interval opens an additional
longitudinal phase space in which new soft modes are generated. The effect of the boost is
represented by a unitary operator $\OM$ acting on the valence state, which describes the
emission, absorption and recombination of the soft quanta. At leading order $\OM$ generates the
Weizs\"acker--Williams gluon field \cite{Kovner:2005pe,Kovner:2006wr}; beyond leading order it
develops multi-particle and colour correlated structures responsible for fluctuations and
correlations \cite{Kovner:2016jfp}. In a companion paper \cite{PaperI}, hereafter referred to as
Paper~I, we constructed the soft gluon wave function and the evolution operator $\OM$ explicitly
up to $\order{g^{2}}$ in pure Yang--Mills theory ($\Nf=0$), fixing all coefficients of the
normal ordered expansion of $\OM$ by unitarity and by direct light-cone perturbation theory
(LCPT) \cite{Lepage:1980fj,Brodsky:1997de,Munier:2025unitary}, and used it to diagonalize the
soft Yang--Mills Hamiltonian through order $g^{2}$.

The purpose of the present work is to include the quark degrees of freedom and thereby obtain
the evolution operator of full QCD to the same accuracy. Quarks enter the problem in two
distinct ways. First, the valence sector of a hadron contains quarks, and in the
Born--Oppenheimer picture \cite{Duan:2024nlr,Duan:2024qev,Duan:2024BOmodBFKL} their only effect on the
soft modes is through their colour charge: the valence charge density acquires a fundamental
representation component, $\rho=\rho_{g}+\rho_{q}$, which obeys the same $SU(N_{c})$ algebra as
in the pure gauge theory. Second, and more importantly, soft gluons can fluctuate into soft
quark--antiquark pairs already at order $g$, soft quarks radiate and absorb soft gluons, and
soft quarks rescatter eikonally on the valence charge. These processes generate new operator
structures in $\OM$. They are the source of the $\Nf$ dependence of NLO small-$x$ evolution
\cite{Balitsky:2006wa,Kovchegov:2006vj,Balitsky:2007feb,Kovner:2013ona,Lublinsky:2016meo}, of the
quark contributions to forward particle production and of the sea quark component of the
projectile wave function, and they must be under control before the full QCD single inclusive
gluon spectrum at NLO \cite{LRS:inprep} can be computed. Light-cone wave functions involving
quarks have been computed in several related contexts (see e.g.
Refs.~\cite{Beuf:2016wdz,Hanninen:2017ddy,Lappi:2016oup,Lappi:2026dglap}); here we organize them into
the unitary operator $\OM$ in the presence of the non-commuting valence charge.

We follow the strategy of Paper~I closely. After the Born--Oppenheimer separation and the
eikonal approximation, we parameterize $\OM$ as a fully normal ordered series in soft gluon,
quark and antiquark creation and annihilation operators with operator valued coefficients that
are functionals of $\rho$. We act with $\OM$ on elementary Fock states --- the vacuum, one, two
and three-gluon states, a single quark (antiquark), a quark--antiquark pair and a
quark--gluon state --- and equate the result, sector by sector, with the corresponding LCWF
components computed in LCPT. This one-to-one matching fixes the coefficients associated with
real transitions, while the unitarity condition $\OM^{\dg}\OM=1$, imposed order by order in the
coupling, determines the remaining ones and provides nontrivial checks of the coefficients that
preserve the particle content. Two results of the construction deserve to be emphasized at the
outset. First, the ansatz that contains only the pair creation and annihilation structures is
not complete: at order $g$ the evolution operator must also contain the coefficients
$\mf F$ and $\mf G$ describing the emission of a soft gluon by a soft quark and by a soft
antiquark, and at order $g^{2}$ the coefficients $\mf P$ and $\bar{\mf P}$ describing the
eikonal rescattering of a soft quark (antiquark) on the valence charge. Without them $\OM$ is
neither unitary nor able to diagonalize the Hamiltonian at order $g$. Second, every second order
coefficient can be written as a sum of products of first order coefficients weighted by ratios
of energy denominators, together with the instantaneous terms; this reduction formula makes all
the unitarity relations transparent and allows a compact presentation of the coefficients with
many external legs.

As an application we use the evolution operator to diagonalize the soft QCD Hamiltonian up to
orders $g$ and $g^{2}$. At order $g$ the Hamiltonian is fully diagonalized once the coefficients
$A$, $C$ of Paper~I are supplemented by $\mf A_{3}$, $\mf F$ and $\mf G$. At order $g^{2}$ all
quark induced off-diagonal terms are removed, the quark loop contribution to the gluon
self-energy and the quark self-energy are scaleless, and the only surviving correction is the
coherent background field energy proportional to $\rho^{2}$. The sole modification relative to
the pure Yang--Mills result is that $\rho$ now contains the valence quark charge.

Our article is structured as follows. In section~\ref{sec:setup} we introduce the light-cone QCD
Hamiltonian with quarks and discuss the Born--Oppenheimer separation and the eikonal
approximation. In section~\ref{sec:lcwf} we review the light-cone wave function formalism and
LCPT and derive the reduction formula for second order amplitudes. In
section~\ref{sec:omega} we construct the perturbative evolution operator with quarks, derive the
unitarity relations and its action on the relevant Fock states, and compute the corresponding
LCWFs. In section~\ref{sec:coeff} we determine all the coefficients. In section~\ref{sec:diag} we
diagonalize the soft QCD Hamiltonian at orders $g$ and $g^{2}$. Section~\ref{sec:concl} contains
the concluding remarks. Technical details are collected in the appendices.
